% ─────────────────────────────────────────────────────────────────────────────
\section{Conclusion}
% ─────────────────────────────────────────────────────────────────────────────

We present a decentralized model of UAV-based search and rescue combining BDI-inspired statechart reasoning with a reversed ACO coordination mechanism. The design deliberately prioritizes interpretability alongside performance: every behavioral mode is a named state, every transition carries an explicit condition, and the coordination mechanism reduces to a single scoring formula. The statechart structure makes the system auditable in a way that a monolithic controller cannot offer. Each state handles exactly one responsibility and the overall organization reflects the MAS design principle of decentralized, structured agency. This makes the system easier to reason about, explain, and debug, even if it does not provide mechanical formula-level tracing from a single parameter change to a specific output difference.

The full-factorial sweep shows that the decentralized controller is not only interpretable but also calibratable. Step length is the strongest performance lever, candidate count introduces a non-monotonic trade-off, and the most fragile configurations cluster in the short-step, high-candidate regime. In contrast, a broad robust region remains stable across seeds and many parameter combinations converge consistently. These results show that local stigmergic rules generate predictable global reliability patterns without centralized coordination, which is precisely the type of emergent behavior this model was designed to expose.

\lipsum[7]

The model's simplifications define the boundaries of its validity. The environment is obstacle-free and two-dimensional. Real terrain constrains UAV trajectories in ways the model does not capture. The sensor is idealized, providing binary detection within a clean geometric cone with no noise, no false positives, and no occlusion. Coordination runs entirely through the shared pheromone grid, which presupposes a lossless shared data structure rather than an explicit radio channel. Victim positions are static. These simplifications are deliberate: they isolate the behavioral mechanisms of interest from confounding physical factors. But they establish the ceiling on the model's direct applicability.

Extending the model toward real-world deployment would require addressing each of these gaps. A probabilistic sensor model would replace the binary threshold with a Bayesian confidence update, accumulating evidence over multiple sensor passes and supporting partial detections. The shared pheromone grid would become a distributed gossip protocol. Each UAV would maintain a local copy and periodically merge it with neighbors within radio range, introducing communication latency and network partitions as new performance variables. Battery management would need hardware-specific discharge curves that account for wind, payload, and temperature. Regulatory constraints (altitude limits, beyond-visual-line-of-sight certification, and mandated collision avoidance margins) would impose hard trajectory constraints beneath the statechart layer.

Despite these gaps, the core design principles validated here transfer cleanly to more physically realistic architectures. Repulsive stigmergy for coverage dispersion, attractive stigmergy for exploitation, and a transparent statechart for inspectable behavioral switching are not artifacts of the simplified model. They are general principles applicable to any representation of recently-searched space, any detection model that produces a confidence signal, and any behavior architecture that supports explicit state modelling. The sensitivity analysis framework developed here provides a methodology that can be carried forward to guide calibration as the model is progressively extended toward deployment.
